
\section{引理~8.7.2 的证明}
\label{sec:proof-lemma-872}

设 $K\subset\RR\times S^{1}$ 是一个\trCompact{}集，并设 $U\supset K$ 是一个\trOpen{}集。那么由\cref{thm:local-regularity-estimate}，存在一个常数 $C_{1}(K,U)$，使得
\[
\|Y\|_{W^{1,p}(K)}\le C_{1}\left(\|\bar{\partial}Y\|_{L^{p}(U)}+\|Y\|_{L^{p}(U)}\right)\text{。}
\]
由于
\[
\|\bar{\partial}Y\|_{L^{p}(U)}\le\|LY\|_{L^{p}(U)}+\|SY\|_{L^{p}(U)}\text{，}
\]
我们便得到
\[
\|Y\|_{W^{1,p}(K)}\le C_{2}\left(\|LY\|_{L^{p}(U)}+\|Y\|_{L^{p}(U)}\right)\text{。}
\]

接下来，让我们假设 $S(s,t)=S(t)$，并在这种情形下证明那条不等式。由上面的结果，我们有
\[
\|Y\|_{W^{1,p}([0,1]\times S^{1})}\le C_{3}\left(\|LY\|_{L^{p}([-1,2]\times S^{1})}+\|Y\|_{L^{p}([-1,2]\times S^{1})}\right)
\]
其中 $C_{3}>0$ 是某个常数。由于 $S$ 不依赖于 $s$，算子 $L$ 在 $s$ 的平移下不变，也就是说
\[
L\bigl(Y(s+s_{0},t)\bigr)=(LY)(s+s_{0},t)\text{。}
\]
由此可知，对每个 $k\in\ZZ$，都有
\[
\|Y\|_{W^{1,p}([k,k+1]\times S^{1})}\le C_{3}\left(\|LY\|_{L^{p}([k-1,k+2]\times S^{1})}+\|Y\|_{L^{p}([k-1,k+2]\times S^{1})}\right)\text{。}
\]
像在\cref{thm:d-bijective-independent-s}的证明中那样，应用不等式 $(a+b)^{p}\le 2^{p}(a^{p}+b^{p})$，我们便得到
\[
\|Y\|^{p}_{W^{1,p}([k,k+1]\times S^{1})}\le 2^{p}C_{3}^{p}\left(\|LY\|^{p}_{L^{p}([k-1,k+2]\times S^{1})}+\|Y\|^{p}_{L^{p}([k-1,k+2]\times S^{1})}\right)\text{。}
\]
随后我们对全体整数 $k$ 取和，就发现\footnote{译注：下式两行中 $\|Y\|$ 项的范数下标，原书均印作 $L^{p}_{\times S^{1}}$（缺 $\mathbf{R}$）；按同式内 $\|LY\|$ 项的下标 $L^{p}(\mathbf{R}\times S^{1})$，疑为漏排，疑应为 $L^{p}(\mathbf{R}\times S^{1})$。此处径改。——译者注}
\[
\begin{aligned}
\|Y\|^{p}_{W^{1,p}(\RR\times S^{1})}&\le C_{4}\left(\|LY\|^{p}_{L^{p}(\RR\times S^{1})}+\|Y\|^{p}_{L^{p}(\RR\times S^{1})}\right)\\
&\le C_{4}\left(\|LY\|_{L^{p}(\RR\times S^{1})}+\|Y\|_{L^{p}(\RR\times S^{1})}\right)^{p}
\end{aligned}
\]
因而当 $S$ 不依赖于 $s$ 时，所求的那条不等式\footnote{译注：原书此处印作 “the desired equality”（所求的等式）；所示实为引理的那条不等式，疑应为 “the desired inequality”（\cref{sec:why-how}中转述引理~8.7.2 处的同款措辞已按同一原则处理）。此处径改。——译者注}便成立。让我们回到一般的情形。设 $D^{\pm}$ 是分别与 $S^{\pm}(t)$ 相联系的受\trPerturbation{}的\trCauchyRiemann{}算子。设 $\varepsilon>0$。若 $M>0$ 充分大，则
\[
\|S(s,t)-S^{+}(t)\|<\varepsilon\qquad\text{对 }s>M\text{ 与 }t\in S^{1}\text{，}
\]
且
\[
\|S(s,t)-S^{-}(t)\|<\varepsilon\qquad\text{对 }s<-M\text{ 与 }t\in S^{1}\text{。}
\]
因而，在 $W^{1,p}\bigl([M,+\infty[\,\times S^{1};\RR^{2n}\bigr)$ 上，我们有 $\|L-D^{+}\|^{\mathrm{op}}<\varepsilon$。

设 $Y\in W^{1,p}(\RR\times S^{1})$，其\trSupport{}含于 $\{s>M\}$ 内。\footnote{译注：自此句起，原书恒以 $Y\in W^{1,p}(\mathbf{R}\times S^{1})$ 为前提，故全证明实际交付的只是引理的那条\trApriori{}估计；而\cref{sec:why-how}对引理~8.7.2 的转述在此之外还多出一句加强断言：“更一般地，若 $Y$ 与 $LY$ 都属于 $L^{p}$，则 $Y\in W^{1,p}$”（原书 p.~455）。要从这一较弱的前提出发，须另行补三步：先由\cref{thm:local-regularity-estimate}得局部的\trRegularity{} $Y\in W^{1,p}_{\mathrm{loc}}$（此时 $\bar{\partial}Y\in L^{p}$：$LY$ 与\trBounded{}的 $SY$ 皆属于 $L^{p}$）；再对\trPlateauFunction{}截断 $\beta Y$ 逐段套用本证明的统一估计；末令截断随 $M$ 增大而趋于原 $Y$——书中并无这一极限论证。该结论本身是标准的、可循此补全，但现证明确实未处理较弱的前提。照排。——译者注}我们有（照我们刚才对 $S$ 不依赖于 $s$ 的情形所做的）
\[
\begin{aligned}
\|Y\|_{W^{1,p}}&\le C_{5}\left(\|D^{+}Y\|_{L^{p}}+\|Y\|_{L^{p}}\right)\\
&=C_{5}\left(\|D^{+}Y\|_{L^{p}([M,+\infty[\,\times S^{1})}+\|Y\|_{L^{p}(\RR\times S^{1})}\right)\\
&\le C_{5}\left(\|LY\|_{L^{p}([M,+\infty[\,\times S^{1})}+\varepsilon\|Y\|_{W^{1,p}([M,+\infty[\,\times S^{1})}+\|Y\|_{L^{p}(\RR\times S^{1})}\right)\\
&=C_{5}\left(\|LY\|_{L^{p}}+\varepsilon\|Y\|_{W^{1,p}}+\|Y\|_{L^{p}}\right)\text{。}
\end{aligned}
\]
对充分小的 $\varepsilon$（$\varepsilon\le 1/2C_{5}$），引理的那条不等式便对 $Y$ 成立。同理，若 $Y$ 的\trSupport{}含于 $\{s<-M\}$ 内，则同样的结论也成立。

现在我们假设 $Y$ 的\trSupport{}含于 $\{|s|>M\}$ 内。我们把它分解成一个和式 $Y=Y^{+}+Y^{-}$，其中项 $Y^{+}$ 的\trSupport{}含于 $\{s>M\}$ 内，而 $Y^{-}$ 的\trSupport{}含于 $\{s<-M\}$ 内。当然，
\[
\|Y^{\pm}\|_{L^{p}}\le\|Y\|_{L^{p}}\qquad\text{且}\qquad\|LY^{\pm}\|_{L^{p}}\le\|LY\|_{L^{p}}\text{，}
\]
于是我们便有
\[
\begin{aligned}
\|Y\|_{W^{1,p}}&\le\|Y^{+}\|_{W^{1,p}}+\|Y^{-}\|_{W^{1,p}}\\
&\le C_{6}\left(\|LY^{+}\|_{L^{p}}+\|LY^{-}\|_{L^{p}}+\|Y^{+}\|_{L^{p}}+\|Y^{-}\|_{L^{p}}\right)\\
&\le 2C_{6}\left(\|LY\|_{L^{p}}+\|Y\|_{L^{p}}\right)\text{。}
\end{aligned}
\]

接下来，我们考虑\trSupport{}含于 $\{-M-1\le s\le M+1\}$ 内的 $Y$ 的情形。证明开头处所用的论证给出一个（依赖于 $M$ 的）常数 $C$，使得
\[
\begin{aligned}
\|Y\|_{W^{1,p}}&=\|Y\|_{W^{1,p}([-M-1,M+1]\times S^{1})}\\
&\le C_{7}\left(\|LY\|_{L^{p}}+\|Y\|_{L^{p}}\right)\text{。}
\end{aligned}
\]

为了收尾，我们使用一个\trPlateauFunction{} $\beta\colon\RR\longrightarrow[0,1]$，它在 $|s|\ge M+1$ 时为零\footnote{译注：原书此处印作 “zero for $|s|\le M+1$ and is $1$ for $|s|\le M$”（在 $|s|\le M+1$ 时为零、在 $|s|\le M$ 时为一）；如此则零集 $\{|s|\le M+1\}$ 将包含“一”集 $\{|s|\le M\}$，自相矛盾，且与下文把两块 $\beta Y$ 与 $(1-\beta)Y$ 分别归入已证诸情形的用法不接；疑应为“在 $|s|\ge M+1$ 时为零、在 $|s|\le M$ 时为一”（只须改第一个不等号的方向）。此处径改。——译者注}，而在 $|s|\le M$ 时为一。我们写出 $Y=\beta Y+(1-\beta)Y$，使得这两块 $\beta Y$ 与 $(1-\beta)Y$ 都满足引理的那条不等式。注意
\[
\begin{aligned}
\|L(\beta Y)\|_{L^{p}}&=\|\beta(s)LY+\beta'(s)Y\|_{L^{p}}\\
&\le\|LY\|_{L^{p}}+K\|Y\|_{L^{p}}
\end{aligned}
\]
对 $L\bigl((1-\beta)Y\bigr)$ 亦然，其中 $K=\sup_{s\in\RR}|\beta'(s)|$。但是这样一来，
\[
\begin{aligned}
\|Y\|_{W^{1,p}}&\le\|\beta Y\|_{W^{1,p}}+\|(1-\beta)Y\|_{W^{1,p}}\\
&\le 2C_{7}\left(\|LY\|_{L^{p}}+K\|Y\|_{L^{p}}\right)\\
&\le C\left(\|LY\|_{L^{p}}+\|Y\|_{L^{p}}\right)\text{，}
\end{aligned}
\]
这就结束了引理的证明。\noqed
